% Session 2 -- from lattice to Bloch Hamiltonian: the recipe.
% Placement: sec. "Multi-orbital band structures", between eqs. (2.12) and (2.13).
%
% NOTE: the intracell hopping t_AB sits OFF the diagonal here -- that is the whole
% point of the annotation, and the reason the adjacent prose was corrected.
%
% The \foreach loop variable is \et, NOT \eta: naming it \eta silently substitutes
% into every "$\eta$" in the panel titles and prints the numeric value instead.
\begin{aqmfigure}{fig:s02:recipe}{Reading a Bloch Hamiltonian off a lattice. The number of orbitals
per unit cell sets the matrix size; an intracell hopping enters the entry its orbital pair indexes
with no phase, an intercell hopping enters the corresponding entry multiplied by $e^{\pm i\varphi}$.
Note that the intracell $A$--$B$ hopping lands \emph{off} the diagonal. For the alternating-hopping
chain (2.13) this gives the two bands $\epsilon_\pm=\pm|1+\eta e^{-i\varphi}|$: gapped for
$\eta\neq1$, and gap-closing exactly at $\eta=1$ -- the Su-Schrieffer-Heeger chain the topology
session picks up.}

\begin{tikzpicture}[x=1cm, y=1cm]

\def\dx{1.55}   % intracell A-B separation
\def\cw{2.65}   % unit-cell period

% ======================= left: the two-orbital chain ========================
\begin{scope}[shift={(0,0)}]
  \node[aqmpanel, anchor=south west] at (-0.75,1.55) {real space};

  % three unit cells
  \foreach \n/\lab in {0/{$n-1$}, 1/{$n$}, 2/{$n+1$}} {
    \begin{scope}[shift={(\n*\cw,0)}]
      \draw[aqmcell] (-0.46,-0.52) rectangle ({\dx+0.55},0.52);
      \node[aqmlabmuted, anchor=north] at ({\dx/2},-1.15) {\lab};
      % strong intracell bond A-B
      \draw[aqmbondstrong, draw=aqmaccent] (0,0) -- (\dx,0);
      \node[aqmsite, minimum size=7pt] at (0,0) {};
      \node[aqmsitealt, minimum size=7pt] at (\dx,0) {};
      \node[aqmlabsmall, anchor=south] at (0,0.12) {$A$};
      \node[aqmlabsmall, anchor=south] at (\dx,0.12) {$B$};
    \end{scope}
  }
  % weak intercell bonds B(n) - A(n+1)
  \foreach \n in {0,1} {
    \draw[aqmbondweak, draw=aqmbroken] ({\n*\cw+\dx},0) -- ({(\n+1)*\cw},0);
  }
  \draw[aqmbondweak, draw=aqmbroken] (-0.68,0) -- (0,0);
  \draw[aqmbondweak, draw=aqmbroken] ({2*\cw+\dx},0) -- ({2*\cw+\dx+0.68},0);

  \node[aqmlabsmall, anchor=south, text=aqmaccent] at ({\cw+\dx/2},0.66)
    {intracell $t_{AB}=1$};
  \node[aqmlabsmall, anchor=north, text=aqmbroken] at ({\cw+\dx+(\cw-\dx)/2},-0.60)
    {intercell $\gamma=\eta$};
\end{scope}

% ======================= right: the Bloch matrix ============================
\begin{scope}[shift={(8.95,0)}]
  \node[aqmpanel, anchor=south] at (1.45,1.55) {Bloch Hamiltonian $H(\varphi)$};

  \node[aqmbox, text width=3.9cm, minimum height=1.5cm] (mat) at (1.45,0)
    {$\displaystyle
      \begin{pmatrix}
        0 & \textcolor{aqmaccent}{1}+\textcolor{aqmbroken}{\eta e^{-i\varphi}}\\[2pt]
        \textcolor{aqmaccent}{1}+\textcolor{aqmbroken}{\eta e^{i\varphi}} & 0
      \end{pmatrix}$};

  \node[aqmlabsmall, anchor=north, align=center, text width=5.0cm] at (1.45,-0.95)
    {\textcolor{aqmaccent}{intracell: no phase}\quad
     \textcolor{aqmbroken}{intercell: $\times\,e^{\pm i\varphi}$}\\[1pt]
     \textcolor{aqmmuted}{matrix size $=$ orbitals per cell $=2$}};
\end{scope}

\draw[aqmarrow] (7.45,0) -- (8.55,0);

% ======================= bottom: the three band panels ======================
\begin{scope}[shift={(0,-6.55)}]
\foreach \px/\et/\ttl/\note in {0/0.45/{$\eta<1$: gapped}/{},
                                4.75/1.0/{$\eta=1$: gap closes}/{gap closes at $\varphi=\pm\pi$},
                                9.50/1.8/{$\eta>1$: gapped again}/{}} {
  \begin{scope}[shift={(\px,0)}]
    \begin{axis}[
      aqmplot,
      width=5.3cm, height=4.3cm,
      title={\ttl},
      xlabel={$\varphi$}, ylabel={$\epsilon_\pm$},
      ylabel style={at={(ticklabel* cs:1.0)}, anchor=south east},
      xmin=-3.9, xmax=3.9, ymin=-3.3, ymax=3.3,
      xtick={-3.14159,0,3.14159}, xticklabels={$-\pi$,$0$,$\pi$},
      ytick=\empty,
      axis lines=middle,
    ]
    % eps_pm = +- |1 + eta e^{-i phi}| = +- sqrt(1 + eta^2 + 2 eta cos phi)
    \addplot[aqmcurve, domain=-3.14159:3.14159, samples=220]
      {sqrt(1 + \et*\et + 2*\et*cos(deg(x)))};
    \addplot[aqmcurve, domain=-3.14159:3.14159, samples=220]
      {-sqrt(1 + \et*\et + 2*\et*cos(deg(x)))};
    % per-panel note, positioned in axis coordinates so it lands correctly
    \node[aqmlabsmall, anchor=north, text=aqmbroken] at (axis cs:0,-2.42) {\note};
    \end{axis}
  \end{scope}
}
\end{scope}

\end{tikzpicture}

\end{aqmfigure}
